\subsection{কলেজের ভর্তি, ক্লাস ও Result-এর প্রয়োজন}
\label{ch2:c10-14}
\begin{examplebox}{একটি বাস্তব সমস্যা থেকে সেবা নির্বাচন}
কলেজে ভর্তি চলাকালে ওয়েবসাইটে ব্যবহারকারীর চাপ বাড়ে। শিক্ষকরা নোট ও অ্যাসাইনমেন্ট ভাগাভাগি করেন। ফলাফলে প্রবেশাধিকার সীমিত রাখা প্রয়োজন। নোট, ই-মেইল, অনলাইন ডকুমেন্ট, উপস্থিতি ও সভার জন্য প্রস্তুত SaaS ব্যবহার করা যায়। নিজস্ব ভর্তি অ্যাপ্লিকেশন তৈরিতে পরিচালিত প্ল্যাটফর্ম বা PaaS সুবিধাজনক। নির্দিষ্ট OS ও সার্ভারের নিয়ন্ত্রণ প্রয়োজন হলে IaaS বিবেচনা করা যায়।

চাহিদা বাড়লে সম্পদ বাড়ানো, পরে কমানো এবং ব্যবহার ও বিল পর্যবেক্ষণ অংশ \ref{ch2:c10-02}-এর বৈশিষ্ট্যের প্রয়োগ। ফলাফলের জন্য প্রতিষ্ঠানের নিজস্ব Private ব্যবস্থা বিবেচনা করা যায়। Public Admission Cloud-এর সঙ্গে অনুমোদিত ডেটা বিনিময়ের সমন্বিত সংযোগ থাকলে Hybrid হতে পারে। গবেষণাগারের সাময়িক বড় প্রক্রিয়াকরণের কাজেও চাহিদামতো সম্পদ নেওয়া উপকারী। বাস্তব নকশায় উপযুক্ত অনুমতি, ব্যাকআপ ও দায়িত্বের সীমা রাখতে হয়।
\end{examplebox}

\subsection{MCQ অনুশীলন}

\par
\textbf{MCQ}

\begin{multicols}{2}\begin{enumerate}
\item\chTwoMCQBlock{Public Cloud সাধারণত কোন মাধ্যমে ব্যবহার করা হয়?
\par ক. শুধু স্থানীয় প্রিন্টারে \quad খ. ইন্টারনেটের মাধ্যমে \par গ. শুধু অফলাইন ডকুমেন্টে \quad ঘ. কোনো নেটওয়ার্ক ছাড়াই }
\item\chTwoMCQBlock{Google Drive কোন ধরনের সেবার উদাহরণ?
\par ক. ক্লাউড স্টোরেজ \quad খ. স্থানীয় ডিস্কের বিন্যাস \par গ. স্থানীয় ফাইল ডিফ্র্যাগমেন্টেশন \quad ঘ. কিবোর্ডের বিন্যাস }
\item\chTwoMCQBlock{SaaS-এর পূর্ণরূপ কী?
\par ক. Storage as a Service \quad খ. Software as a Service \par গ. System as a Switch \quad ঘ. Server and Scanner }
\item\chTwoMCQBlock{IaaS প্রধানত কী দেয়?
\par ক. প্রস্তুত সফটওয়্যার \quad খ. কম্পিউটিং অবকাঠামো \par গ. কিবোর্ডের বিন্যাস \quad ঘ. মনিটরের উজ্জ্বলতা }
\item\chTwoMCQBlock{অ্যাপ্লিকেশন তৈরি ও চালাতে PaaS প্রধানত কাদের জন্য উপযোগী?
\par ক. সফটওয়্যার উন্নয়নকারী \quad খ. শুধু টাইপিস্ট \par গ. চিত্রশিল্পী \quad ঘ. বিদ্যুৎমিস্ত্রি }
\item\chTwoMCQBlock{Public Cloud সাধারণত কারা পরিচালনা করেন?
\par ক. তৃতীয় পক্ষের সেবাদাতা \quad খ. শুধু একজন শিক্ষার্থী \par গ. একটি স্থানীয় প্রিন্টার \quad ঘ. একটি মাউস }
\item\chTwoMCQBlock{Public ও Private Cloud অবকাঠামো সমন্বিতভাবে যুক্ত হলে কোন মডেল হয়?
\par ক. Hybrid Cloud \quad খ. শুধু Public Cloud \par গ. শুধু Private Cloud \quad ঘ. শুধু Community Cloud }
\item\chTwoMCQBlock{ক্লাউড কম্পিউটিংয়ের একটি সুবিধা কোনটি?
\par ক. প্রয়োজনমতো সম্পদ বাড়ানো বা কমানো \quad খ. শুধু অফলাইনে ব্যবহার \par গ. কেবল পেঁচানো \quad ঘ. মনিটর পরিষ্কার করা }
\item\chTwoMCQBlock{ক্লাউড ব্যবহারে কোন ঝুঁকি বিবেচনা করতে হয়?
\par ক. অননুমোদিত প্রবেশের ঝুঁকি \quad খ. প্রিন্টারের রং \par গ. কিবোর্ডের আকার \quad ঘ. মাউস প্যাড }
\item\chTwoMCQBlock{ক্লাউড অ্যাকাউন্টের নিরাপত্তা বাড়াতে কোন ব্যবস্থা উপযোগী?
\par ক. দুর্বল পাসওয়ার্ড \quad খ. MFA \par গ. পাসওয়ার্ড ভাগাভাগি \quad ঘ. সবার জন্য লগইন উন্মুক্ত রাখা }
\end{enumerate}\end{multicols}

\textbf{উত্তরমালা:}
১. খ, ২. ক, ৩. খ, ৪. খ, ৫. ক, ৬. ক, ৭. ক, ৮. ক, ৯. ক, ১০. খ

\smallskip
